% !PRINTED 510
\weberpage{510}
%<<CONT
sein. Den Inbegriff $R$ aller gemeinsamen Permutationen von $Q$
und $Q'$ nennen wir den \gesperrt{grössten gemeinschaftlichen Theiler},
oder auch den \gesperrt{Durchschnitt} von $Q$ und $Q'$. Dies $R$ ist immer eine
Gruppe, denn wenn $\pi_1$ und $\pi_2$ beide sowohl in $Q$ als in $Q'$
vorkommen, so muss auch $\pi_1 \pi_2$ in $Q$ und in $Q'$, also auch in $R$
vorkommen. Der Begriff ist sofort übertragbar auf mehrere
Gruppen $Q, Q', Q'' \ldots$

Ist $\psi$ eine zu $Q$ und $\psi'$ eine zu $Q'$ gehörige Function, so
können wir die rationalen Zahlen $x, x'$ (nach §. 143, 1.) so bestimmen,
dass $\omega = x\psi + x'\psi'$ eine zu $R$ gehörige Function
wird. Denn jedenfalls gestattet $\omega$ die Permutationen von $R$, da
$\psi$ und $\psi'$ sie gestatten. Ist dann $\pi$ eine nicht in $R$ enthaltene
Permutation aus $P$, so ist sicher nicht zugleich $\psi = \psi(\pi)$ und
$\psi' = \psi'(\pi)$; also können wir $x, x'$ so bestimmen, dass auch nicht
$\omega = \omega(\pi)$ wird. Wenn wir also gleichzeitig $\psi$ und $\psi'$ und folglich
 $\omega$ adjungiren, so reducirt sich die Gruppe des Körpers $N$
nach 3. auf $R$. Ebenso können wir bei mehr als zwei Gruppen $Q$
schliessen, und erhalten den Satz:

4. \gesperrt{Sind $Q, Q', Q'' \ldots$ Theiler von $P$, $R$ ihr Durchschnitt,
$\psi, \psi', \psi'' \ldots$ Functionen, die zu $Q, Q', Q'' \ldots$
gehören, so reducirt sich die Gruppe des
Körpers $N$ durch gleichzeitige Adjunction von
$\psi, \psi', \psi'' \ldots$ auf $R$.}

Wenn wir nicht bloss eine Wurzel $\psi$ der Hülfsgleichung
$\varphi(t) = 0$ adjungiren, sondern alle mit $\psi$ conjugirten Grössen
$\psi, \psi_1, \psi_2 \ldots \psi_{j-1}$, wenn wir also aus $\Omega$ den Körper

\[
\Omega'' = \Omega(\psi, \psi_1, \psi_2 \ldots \psi_{j-1})
\]

ableiten, so ist nach diesem Satze der Erfolg der, dass die
Gruppe von $N$ in $\Omega''$ der grösste gemeinschaftliche Theiler $R$
aller mit $Q$ conjugirten Gruppen wird, so dass (nach §. 154)
$R$ der grösste gemeinschaftliche Theiler aller Gruppen $\pi^{-1}Q\pi$
ist, wenn $\pi$ die Permutationen von $P$ durchläuft.

\gesperrt{Von besonderem Interesse ist nun der Fall, dass alle conjugirten
Gruppen $\pi^{-1}Q\pi$ mit einander identisch sind.} In diesem
Falle ist nach dem Theorem 2. jede der conjugirten Grössen
$\psi, \psi_1, \psi_2 \ldots \psi_{j-1}$ in $\Omega(\psi)$ enthalten, die Körper
$\Omega(\psi), \Omega(\psi_1) \ldots \Omega(\psi_{j-1})$ und $\Omega(\psi, \psi_1 \ldots \psi_{j-1})$ sind identisch,
und $\Omega(\psi)$ ist ein \gesperrt{Normalkörper} über $\Omega$.
